\scanpage{6}
\begin{center}- 42 -\end{center}
\hfill n$^\circ$ 387
\begin{center}§ 5.\end{center}

1) Soient $A$ un anneau intègre, $K$ son corps des fractions, $T$ une topologie linéaire sur $A$ (chap. III, § 2, exerc. 21).

a) Pour que les voisinages de $0$ pour $T$ constituent un système fondamental de voisinages de $0$ pour une topologie $T_K$ compatible avec la structure d'anneau de $K$, il faut et il suffit que $T$ soit la topologie $T_u(A)$ (chap. III, § 2, exerc. 24) ; $A$ est alors une partie bornée pour $T_K$ et $T_K$ est une topologie séparée localement bornée (\uline{Top. gén.}, chap. III, \unsure{3|2}$^{\mathrm e}$ éd., § 6, exerc. 12 et 20e). Pour que $K$ soit complet (resp. linéairement compact, resp. strictement linéairement compact (chap. III, § 2, exerc. 21)) pour $T_K$, il faut et il suffit que $A$ le soit pour $T$.

b) Pour que la topologie $T_K$ (où $T=T_u(A)$) soit compatible avec la structure de corps de $K$, il faut et il suffit que le radical $\underline R(A)$ de $A$ soit $\ne0$.

\textsuperscript{*}2) Soient $K$ un corps (commutatif), $T$ une topologie séparée sur $K$, compatible avec la structure d'anneau de $K$. Pour que la topologie $T$ soit définie par une valuation de $K$ ou une valeur absolue sur $K$, il faut et il suffit que $T$ soit localement rétrobornée (\uline{Top. gén.}, chap. III, 3$^{\mathrm e}$ éd., § 6, exerc. 22. S'il existe dans $K$ des éléments topologiquement nilpotents utiliser l'exerc. 22 d) de \uline{Top. gén.}, chap. III, 3$^{\mathrm e}$ éd., § 6 et l'exerc. 13 de \uline{Top. gén.}, chap. IX, 2$^{\mathrm e}$ éd., § 3. Dans le cas contraire, utiliser l'exerc. 22 f) de \uline{Top. gén.}, chap. III, 3$^{\mathrm e}$ éd., § 6). *

\textsuperscript{*}3) Soient $A$ un anneau noethérien intègre, $K$ son corps des fractions, $T$ la topologie $T_u(A)$ sur $A$, $T_K$ la topologie correspondante sur $K$ (exerc. 1).

a) Si $A$ est un anneau de Zariski de radical $\underline r\ne0$ et si $A$ est complet pour la topologie $\underline r$-adique, montrer que $A$ est complet pour la topologie $T_u$ (\uline{Top. gén.}, chap. III, 3$^{\mathrm e}$ éd., § 3, n$^\circ$ 5, cor. 2 de la prop. 9).

b) On suppose que $T_K$ est non discrète et est définie par une valuation $v$ de $K$. Montrer alors que $v$ est une valuation discrète et que $A$ est un sous-anneau ouvert de l'anneau $V$ de $v$ ; réciproque. (En utilisant l'hypothèse, qui entraîne $A\ne K$, et la prop. 4 du § 4, n$^\circ$ 1, on peut supposer que $A\subset V$. En utilisant
\mspage{5}
